\documentclass[11pt]{article}
\usepackage[T1]{fontenc}
\usepackage[margin=1in]{geometry}
\usepackage{amsmath,amssymb,amsthm,hyperref}
\setlength{\emergencystretch}{3em}
\newtheorem{theorem}{Theorem}
\newtheorem{lemma}[theorem]{Lemma}
\newcommand{\R}{\mathbb R}
\newcommand{\abs}[1]{\lvert#1\rvert}
\newcommand{\norm}[1]{\lVert#1\rVert_1}
\title{A minimum-side refinement of the taxicab\\Erd\H{o}s--Mordell inequality}
\author{Myeongjin Kang}
\date{26 October 2026 --- research draft}
\begin{document}
\maketitle
\begin{abstract}
For an arbitrary taxicab triangle and a point in its closed convex hull, let $R_A,R_B,R_C$ be the vertex distances and $r_a,r_b,r_c$ the supporting-line distances. We prove
\[
 R_A+R_B+R_C\ge\frac32\bigl(r_a+r_b+r_c+\min(r_a,r_b,r_c)\bigr).
\]
An explicit strictly interior configuration attains equality. This strengthens the unrestricted three-halves conjecture recorded as MathDB 341671. The proof reduces to a polynomial inequality, a threshold argument and two square identities. A complete Lean certificate covers the refinement, the ordinary triangle interior and both optimal coefficients. Final novelty comparison remains in progress. This draft is not yet counted as an accepted new discovery.
\end{abstract}
\noindent\textit{Provenance.} AI-assisted research. Petrovi\'c, Male\v{s}evi\'c and Banjac posed the unrestricted three-halves conjecture. The old coefficient's sharpness is known prior work. The candidate contribution here is the additive minimum-side refinement, treated as one proposed result with the source conjecture as a consequence. Original prose is CC BY 4.0 and original code is MIT.

\section{Exact geometric statement}
For $X=(x_1,x_2)\in\R^2$, put $\norm X=\abs{x_1}+\abs{x_2}$. The distance to a line $\ell:\alpha x+\beta y+\gamma=0$ is
\begin{equation}\label{eq:distance}
 d_1(P,\ell)=\frac{\abs{\alpha P_1+\beta P_2+\gamma}}{\max(\abs\alpha,\abs\beta)},\qquad(\alpha,\beta)\ne(0,0).
\end{equation}
This is the exact supporting-line convention in the primary source~\cite{pmb}, equations (5)--(6). No change of metric, rotation or coordinate normalization is used.

For distinct vertices $A,B,C$ and $P\in\operatorname{conv}\{A,B,C\}$, let $R_A=\norm{A-P}$, cyclically, and let $r_a=d_1(P,BC)$, cyclically. Each side symbol denotes its full supporting line. Write $m=\min(r_a,r_b,r_c)$.
\begin{theorem}\label{thm:main}
For all such configurations,
\begin{equation}\label{eq:main}
 2(R_A+R_B+R_C)\ge3(r_a+r_b+r_c+m).
\end{equation}
The factor $3/2$ in the equivalent sum-of-vertex-distances inequality cannot be increased, even for strictly interior points of nondegenerate triangles. With that factor fixed, the coefficient one of $m$ cannot be increased either.
\end{theorem}
The original three-halves inequality follows because $m\ge0$. It is a consequence of the same result, not a separately counted discovery.

\section{A polynomial scalar bound}
For six nonnegative numbers define
\[
 H(A,B,C;u,v,w)=ABC-Avw-Buw-Cuv-2uvw.
\]
Paired relabelling of $(A,u),(B,v),(C,w)$ leaves $H$ unchanged.
\begin{lemma}\label{lem:threshold}
Suppose
\[
 u\le\min(B,C),\qquad v\le\min(A,C),\qquad w\le\min(A,B),
\]
and $2(A+B+C)=3(u+v+w+\min(u,v,w))$. Then $H\le0$.
\end{lemma}
\begin{proof}
Relabel so that $0\le u\le v\le w$. Set
\[
 y=w-v,\quad a=A-w,\quad b=B-w,\quad c=C-v.
\]
These are nonnegative, and the threshold equation becomes
\begin{equation}\label{eq:constraint}
 2(a+b+c)+y=6u.
\end{equation}
Consider the polynomial
\[
 h(t)=H(t+y+a,t+y+b,t+c;u,t,t+y).
\]
It is quadratic, with quadratic coefficient $q=b+c-4u\le-u\le0$. Its derivative at $t=u$ is
\[
 D=u(b+c-8u-4y)+b(y+a)+c(2y+a+b).
\]
Using~\eqref{eq:constraint}, direct expansion gives the exact identity
\begin{align}
12D={}&-31u^2-[3(b+c)-8u]^2-(3y-u)^2\notag\\
 &-12yb-3(b-c)^2-18ya\le0.\label{eq:sos}
\end{align}
Since
\[
 h(t)-h(u)=(t-u)\bigl[D+q(t-u)\bigr],
\]
we have $h(v)\le h(u)$. This reduces the problem to equal first two side variables $u,u$ and third side variable $w'=u+y$. The three reduced vertex lengths are $u+y+a,u+y+b,u+c$; they remain nonnegative, satisfy the caps and obey the corresponding threshold equation.

In this equal-side case replace the first two lengths by their mean $s$. This increases $H$ by $C(A-B)^2/4\ge0$. Dropping the prime on the new third side variable, the remaining conditions are
\[
 w\ge u\ge0,\quad s\ge w,\quad C\ge u,\quad4s+2C=9u+3w.
\]
The polynomial factors as
\[
 H(s,s,C;u,u,w)=(s+u)\bigl[C(s-u)-2uw\bigr].
\]
Thus it is enough to show $C(s-u)\le2uw$.

If $3w\le8u$, expansion under the displayed linear constraint gives
\[
 32\bigl[2uw-C(s-u)\bigr]=(4C-5u-3w)^2+(w-u)(25u-9w)\ge0,
\]
because $w\ge u$ and $25u-9w\ge u\ge0$. If $3w\ge8u$, put $d=s-w\ge0$. Then
\[
 2\bigl[2uw-C(s-u)\bigr]=(w-3u)^2+d(5w-13u)+4d^2\ge0,
\]
because $5w-13u\ge u/3\ge0$. These two cases prove the threshold assertion, including all zero-variable boundary cases.
\end{proof}

\begin{lemma}\label{lem:cubic}
For nonnegative $A,B,C,u,v,w$ satisfying the same caps, if $H\ge0$, then
\[
 2(A+B+C)\ge3\bigl(u+v+w+\min(u,v,w)\bigr).
\]
\end{lemma}
\begin{proof}
Suppose the conclusion failed, and put
\[
 S=A+B+C,\qquad T=\tfrac32\bigl(u+v+w+\min(u,v,w)\bigr),\qquad\delta=(T-S)/3>0.
\]
Increasing all three vertex variables by $\delta$ preserves the caps and reaches the threshold. The exact change is
\begin{align*}
&H(A+\delta,B+\delta,C+\delta;u,v,w)-H(A,B,C;u,v,w)\\
&\quad=\delta\bigl[(AB-uv)+(AC-uw)+(BC-vw)\bigr]
 +\delta^2(A+B+C)+\delta^3>0.
\end{align*}
Each bracketed difference is nonnegative by the caps, and $\delta^3>0$. The increased value of $H$ is therefore strictly positive, contradicting Lemma~\ref{lem:threshold}.
\end{proof}

\begin{lemma}[Barycentric scalar form]\label{lem:weighted}
Let all six variables be nonnegative and satisfy the same caps. Suppose $a,b,c\ge0$, $a+b+c=1$, and
\[
 (1-a)u\le aA,\qquad(1-b)v\le bB,\qquad(1-c)w\le cC.
\]
Then the conclusion of Lemma~\ref{lem:cubic} holds.
\end{lemma}
\begin{proof}
If $A,B,C>0$, the weighted inequalities imply
\[
 \frac{u}{A+u}+\frac{v}{B+v}+\frac{w}{C+w}\le a+b+c=1.
\]
Clearing the three positive denominators gives exactly $H\ge0$, so apply Lemma~\ref{lem:cubic}.

In the general case replace the three vertex variables by $A+\varepsilon,B+\varepsilon,C+\varepsilon$, with $\varepsilon>0$. The weighted inequalities and caps remain valid. The positive-length result gives the desired inequality with an additional $6\varepsilon$ on the vertex side. If the desired conclusion failed, choose $\varepsilon$ equal to one twelfth of its positive gap to obtain a contradiction. No limit theorem or numerical approximation is used.
\end{proof}

\section{The exact geometric bridge}
Choose nonnegative barycentric coordinates $P=aA+bB+cC$ with $a+b+c=1$. For any line function $\ell(X)=\alpha X_1+\beta X_2+\gamma$ vanishing on $B,C$, the inequality
\[
 \abs{\alpha x+\beta y}\le\max(\abs\alpha,\abs\beta)(\abs x+\abs y)
\]
implies $r_a\le R_B,R_C$. The affine identity gives $\ell(P)=a\ell(A)$, and hence
\[
 (1-a)\ell(P)=a\bigl(\ell(A)-\ell(P)\bigr).
\]
Since $a,1-a\ge0$, take absolute values, apply the same coordinate bound and divide by the positive line scale. This proves $(1-a)r_a\le aR_A$. Cyclic permutation gives the other five comparisons. Lemma~\ref{lem:weighted}, applied to $(R_A,R_B,R_C;r_a,r_b,r_c)$, proves Theorem~\ref{thm:main}'s inequality over the full closed triangle. Its restriction to the ordinary topological interior requires no additional hypothesis.

\section{Equality and optimal coefficients}
Take
\[
 P=(0,0),\quad A=(3,0),\quad B=(0,3),\quad C=(-3/2,-3/2).
\]
The barycentric coordinates of $P$ are $(1/4,1/4,1/2)$, so it lies strictly inside this nondegenerate triangle. Direct substitution in~\eqref{eq:distance} gives
\[
 R_A=R_B=R_C=3,\qquad(r_a,r_b,r_c)=(1,1,3),\qquad m=1.
\]
Thus both sides of~\eqref{eq:main} equal18. Increasing either the outer factor $3/2$ or the coefficient one of $m$ violates the inequality at this fixed strictly interior configuration. This equality case is part of the same candidate refinement; it receives no separate discovery credit.

\section{Formal verification and novelty status}
The companion development uses explicit coordinate $L^1$ distances, not the default maximum norm on a product of two real lines. All 68 local declarations in the six-module certificate pass fresh isolated replay against pinned Lean 4.34.1 and mathlib. A separate independent replay also passed, with only the standard axioms \texttt{propext}, \texttt{Classical.choice} and \texttt{Quot.sound}. The verified scope includes the generic nonnegative cubic criterion, the weighted scalar inequality, all side-variable orderings and zero cases, the ordinary closed convex hull and its topological interior, the new equality configuration and both optimal-coefficient statements. Independent mathematical and source reviews found no defect. No generic normed-space theorem is claimed as a separate formal result.

The source~\cite{pmb} proves the old lower bound for specified orientations and leaves the arbitrary-triangle assertion as Conjecture 1. The optimality of the old coefficient $3/2$ is not claimed new: Kurusa~\cite[Section 7.3.4]{kurusa} attributes a configuration of that ratio to Ghandehari--Martini~\cite[Figure 3]{gm}. Our proposed contribution is the stronger universal inequality containing the additional minimum-side term.

Two independent mathematical checks found no flaw in the refinement's proof. Searches have not yet located the same additive minimum-side inequality, but unsuccessful searches do not establish absolute priority. The closest 2018 paper has not yet been obtained in full from an accessible primary source. The final prior-work comparison remains required before strict discovery acceptance; independent full-certificate review has passed. No separate count is assigned to the old conjecture, the scalar lemmas, an equality example, a boundary case, or a normed-space reformulation.

\clearpage
\begin{thebibliography}{9}
\bibitem{pmb} M. Petrovi\'c, B. Male\v{s}evi\'c and B. Banjac, \emph{On the Erd\H{o}s--Mordell inequality for triangles in taxicab geometry}, Journal of Mathematical Inequalities 14(4) (2020), 1299--1319. \href{https://doi.org/10.7153/jmi-2020-14-84}{doi:10.7153/jmi-2020-14-84}. Preprint \href{https://arxiv.org/abs/1901.08758}{arXiv:1901.08758}, Conjecture 1.
\bibitem{gm} M. Ghandehari and H. Martini, \emph{On the Erd\H{o}s--Mordell inequality for normed planes and spaces}, Studia Scientiarum Mathematicarum Hungarica 55(2) (2018), 174--189. \href{https://doi.org/10.1556/012.2018.55.2.1392}{doi:10.1556/012.2018.55.2.1392}.
\bibitem{kurusa} \'A. Kurusa, \emph{Metric Characterizations of Projective-Metric Spaces}, in A. Papadopoulos (ed.), Surveys in Geometry II (2024), 225--256, Section 7.3.4. \href{https://doi.org/10.1007/978-3-031-43510-2_7}{doi:10.1007/978-3-031-43510-2\_7}.
\end{thebibliography}
\end{document}
